\documentclass[tikz]{standalone}
\usepackage{amsmath}
\definecolor{bgcol}{HTML}{171D2B}
\definecolor{accent}{HTML}{F7B538}
\definecolor{setcol}{HTML}{7DA2C9}
\pagecolor{bgcol}
\begin{document}

% Gambar 1: Segitiga ABC dengan sudut luar di C (110 derajat)
% A=(0,0), B=(5,0), C=(50:4.61) = (2.96, 3.53)
% Sudut A = 50, sudut B = 60, sudut C = 70 derajat
% D = perpanjangan BC melewati C, arah 120 derajat, 1.5 cm dari C
\begin{tikzpicture}[scale=1.0, color=white, font=\large]
  \coordinate (A) at (0, 0);
  \coordinate (B) at (5, 0);
  \coordinate (C) at (50:4.61);
  \coordinate (D) at (2.21, 4.83);

  % Segitiga ABC
  \draw[thick] (A) -- (B) -- (C) -- cycle;
  % Perpanjangan sisi BC, putus-putus warna emas
  \draw[thick, dashed, accent] (C) -- (D);

  % Busur sudut dalam (biru muda)
  \draw[setcol] (A) ++(0:0.8) arc (0:50:0.8);
  \draw[setcol] (B) ++(120:0.8) arc (120:180:0.8);
  \draw[setcol] (C) ++(230:0.8) arc (230:300:0.8);
  % Busur sudut luar 110 derajat (emas, jari-jari lebih besar)
  \draw[accent] (C) ++(120:1.1) arc (120:230:1.1);

  % Nama titik
  \node[below left] at (A) {$A$};
  \node[below right] at (B) {$B$};
  \node[above] at (C) {$C$};
  \node[above left, accent] at (D) {$D$};

  % Label sudut
  \node[setcol] at (1.04, 0.49) {$50^{\circ}$};
  \node[setcol] at (4.00, 0.58) {$60^{\circ}$};
  \node[setcol] at (2.86, 2.38) {$70^{\circ}$};
  \node[accent] at (1.47, 3.66) {$110^{\circ}$};
\end{tikzpicture}

% Gambar 2: Dua sudut dalam dijajarkan pada garis lurus: 50 + 60 = 110
% O = titik C, garis mendatar = perpanjangan BC, sinar 110 derajat = CA
\begin{tikzpicture}[scale=1.0, color=white, font=\large]
  \coordinate (O) at (0, 0);

  % Garis lurus (perpanjangan BC)
  \draw[thick] (-3.5, 0) -- (3.5, 0);
  % Sinar pembatas sudut A (50 derajat) dan sinar CA (110 derajat)
  \draw[thick] (O) -- (50:3);
  \draw[thick] (O) -- (110:3);

  % Dua sektor: sudut A 50 derajat dan sudut B 60 derajat
  \draw[setcol] (O) ++(0:1.0) arc (0:50:1.0);
  \draw[setcol] (O) ++(50:1.0) arc (50:110:1.0);

  % Nama titik dan ujung garis
  \node[below] at (O) {$C$};
  \node[left] at (-3.5, 0) {$D$};
  \node[right] at (3.5, 0) {$B$};
  \node[above left] at (110:3.3) {$A$};

  % Label sektor dan sudutnya
  \node[setcol] at (1.27, 0.59) {$50^{\circ}$};
  \node[setcol] at (0.24, 1.38) {$60^{\circ}$};
  \node at (2.08, 0.97) {$\angle A$};
  \node at (0.40, 2.27) {$\angle B$};

  % Jumlah kedua sudut = sudut luar
  \node[accent] at (0, -0.9) {$50^{\circ} + 60^{\circ} = 110^{\circ}$};
\end{tikzpicture}

\end{document}
